
\section{Plano lateral}
\label{sec:plano_lateral}

A análise do plano lateral permite estudar o movimento do braço numa vista de lado, utilizando duas coordenadas: $d$, que representa a distância frontal, e $h$, que representa a altura do centro de preensão relativamente à base.

Nesta análise, considera-se $q_1=0^\circ$ como referência, uma vez que a rotação da base não altera as coordenadas $d$ e $h$. Quando J1 roda, o plano lateral acompanha essa rotação, sendo a sua influência considerada posteriormente no plano horizontal.

Como ilustrado na Figura~\ref{fig:lateral}, as juntas J2, J3 e J4 possuem eixos de rotação paralelos e perpendiculares ao plano de movimento. Assim, a posição do braço pode ser calculada através das projeções dos seus elos, utilizando uma abordagem geométrica planar.

\begin{figure}[H]
    \centering
    \resizebox{0.6\textwidth}{!}{%Vista lateral do robô
% (L4 desenhado inclinado para se ver o ângulo theta_C; a garra
%  é perpendicular a L4, por isso alpha = theta_C - 90 graus)
\begin{tikzpicture}[scale=0.3,
                    % Estilos
                        link/.style={line width=2.2pt, line cap=round},
                        junta/.style={circle, draw, fill=white, line width=1pt, inner sep=2.2pt},
                        cota/.style={gray!70!black, thin},
                        ang/.style={->, thin},
                        qang/.style={->, thin, red!75!black},
                        qcota/.style={red!60, thin, dashed}]

  % Coordenadas das juntas
    \coordinate (J1) at (0,0);
    \coordinate (J2) at (0,13.2);
    \coordinate (J3) at (3.8,23.5);
    \coordinate (J4) at (13.2,21.9);
    \coordinate (J5) at (20.06,24.40);
    \coordinate (J6) at (25.26,10.12);

  % Eixos
    \draw[->, thin] (-4,0) -- (31,0) node[below left] {$d$};
    \draw[->, thin] (J1) -- (0,29) node[left] {$h$};

  % Base
    \fill[gray!25] (-3,-0.8) rectangle (3,0);
    \draw[thick] (-5,0) -- (30.5,0);

  % Ângulos absolutos (theta_A, theta_B, theta_C e alpha, a partir da horizontal)
    \draw[cota, dashed] (J2) -- ++(6,0);
    \draw[ang] ($(J2)+(4,0)$) arc[start angle=0, end angle=69.75, radius=4] node[pos=0.55, right=1pt] {$\theta_A$};
    \draw[cota, dashed] (J3) -- ++(9,0);
    \draw[ang] ($(J3)+(8,0)$) arc[start angle=0, end angle=-9.66, radius=8] node[pos=0.5, right=2pt] {$\theta_B$};
    \draw[cota, dashed] (J4) -- ++(7,0);
    \draw[ang] ($(J4)+(4.5,0)$) arc[start angle=0, end angle=20, radius=4.5] node[pos=0.5, right=1pt] {$\theta_C$};
    \draw[cota, dashed] (J5) -- ++(5,0);
    \draw[ang] ($(J5)+(3.5,0)$) arc[start angle=0, end angle=-70, radius=3.5] node[pos=0.45, right=2pt] {$\alpha$};

  % Ângulos das juntas (q2, q3, q4, relativos ao elo anterior)
    \draw[qcota] (J2) -- ++(0,7.5);
    \draw[qang] ($(J2)+(0,6.5)$) arc[start angle=90, end angle=69.75, radius=6.5] node[pos=0.5, left=3pt] {$q_2 < 0$};
    \draw[qcota] (J3) -- ++(69.75:5);
    \draw[qang] ($(J3)+(69.75:4)$) arc[start angle=69.75, end angle=-9.66, radius=4] node[pos=0.35, above right=0pt] {$q_3 < 0$};
    \draw[qcota] (J4) -- ++(-9.66:5);
    \draw[qang] ($(J4)+(-9.66:3)$) arc[start angle=-9.66, end angle=20, radius=3];
    \node[red!75!black] at ($(J4)+(3.6,-3)$) {$q_4 > 0$};

  % Braços
    \draw[link] (J1)  -- node[pos=0.55, right] {$L_1$} (J2);
    \draw[link] (J2) -- node[right] {$L_2$} (J3);
    \draw[link] (J3) -- node[below left=0pt and -2pt] {$L_3$} (J4);
    \draw[link] (J4) -- node[above=1pt] {$L_4$} (J5);
    \draw[link] (J5) -- node[right=2pt] {$D = L_6 + L_7$} (J6);

  % Garra (perpendicular a L4)
    \begin{scope}[shift={(J6)}, rotate=20]
      \draw[line width=1.4pt] (-1.1,1.5) -- (-1.1,-0.5) (1.2,1.5) -- (1.2,-0.5);
      \draw[line width=1.4pt] (-1.1,1.5) -- (1.2,1.5);
    \end{scope}

  % Juntas
    \node[junta] at (J2) {};  \node[left=5pt]  at (J2) {J2 (ombro)};
    \node[junta] at (J3) {};  \node[above left=3pt and 0pt] at (J3) {J3 (cotovelo)};
    \node[junta] at (J4) {};  \node[below=5pt] at (J4) {J4 (pulso)};
    \node[junta] at (J5) {};  \node[above=4pt] at (J5) {J5};
    \fill (J6) circle (6pt);   \node[below right=4pt and 2pt] at (J6) {centro de preensão $(d, h)$};

  % Guias
    \draw[cota, dashed] (J6) -- (25.26,0);
    \draw[cota, dashed] (J6) -- (0,10.12);

\end{tikzpicture}}
    \caption{Vista lateral do manipulador no seu plano vertical.}
    \label{fig:lateral}
\end{figure}

\subsection{Ângulos absolutos}

Os ângulos articulares $q_i$ representam a rotação de cada junta relativamente à sua configuração de referência. No entanto, para calcular as projeções dos elos, é necessário conhecer a orientação de cada um relativamente à horizontal.

Na configuração de referência, quando $q_2=0^\circ$, o elo $L_2$ encontra-se na vertical, formando um ângulo de $90^\circ$ com a horizontal. Por esse motivo, a sua orientação absoluta é obtida somando $90^\circ$ ao ângulo da junta J2.

Os elos seguintes dependem também das rotações das juntas anteriores. Assim, obtêm-se:

\begin{equation}
    \theta_A=q_2+\frac{\pi}{2}, \qquad
    \theta_B=\theta_A+q_3, \qquad
    \theta_C=\theta_B+q_4, \qquad
    \alpha=\theta_C-\frac{\pi}{2}
    \label{eq:angulos}
\end{equation}

Os ângulos $\theta_A$, $\theta_B$ e $\theta_C$ representam as orientações absolutas dos elos $L_2$, $L_3$ e $L_4$, respetivamente.

Como o eixo longitudinal da garra é perpendicular ao elo $L_4$, a sua orientação no plano lateral, representada por $\alpha$, é obtida subtraindo $90^\circ$ a $\theta_C$.

Deste modo, quando $\alpha=-90^\circ$, a garra encontra-se orientada verticalmente para baixo.

\subsection{Posição de cada junta}

Conhecidas as orientações dos elos, é possível calcular a posição de cada junta através das suas projeções.

Para um elo de comprimento $L$ e orientação $\theta$, a projeção horizontal é dada por $L\cos\theta$ e a vertical por $L\sin\theta$.

O cálculo é realizado de forma sucessiva, começando na base e acrescentando a contribuição de cada elo.

A junta J2, correspondente ao ombro, encontra-se a uma altura $L_1$. A partir desta posição, o elo $L_2$ permite determinar J3, o elo $L_3$ permite determinar J4 e o elo $L_4$ permite determinar J5.

Assim, considerando a base como origem do referencial, obtêm-se:

\begin{align}
    \text{J2 (ombro):}\quad
    &(d_2,h_2)=(0,L_1)
    \label{eq:j2}\\
    \text{J3 (cotovelo):}\quad
    &(d_3,h_3)=
    (d_2+L_2\cos\theta_A,\;
    h_2+L_2\sin\theta_A)
    \label{eq:j3}\\
    \text{J4 (pulso):}\quad
    &(d_4,h_4)=
    (d_3+L_3\cos\theta_B,\;
    h_3+L_3\sin\theta_B)
    \label{eq:j4}\\
    \text{J5:}\quad
    &(d_5,h_5)=
    (d_4+L_4\cos\theta_C,\;
    h_4+L_4\sin\theta_C)
    \label{eq:j5}
\end{align}

Estas expressões mostram que a posição de cada junta depende da posição anterior e da orientação do elo que as une.

Em particular, as juntas J2, J3 e J4 formam um sistema equivalente a um manipulador planar de dois graus de liberdade, estudado nas aulas \cite{bicho_2dof}.

Considerando $L_A=L_2$, $L_B=L_3$, $\theta_1=\theta_A$ e $\theta_2=q_3$, a posição de J4 relativamente a J2 pode ser escrita como:

\begin{equation}
    \begin{aligned}
        x_2 &= L_A\cos\theta_1+
        L_B\cos(\theta_1+\theta_2)\\
        y_2 &= L_A\sin\theta_1+
        L_B\sin(\theta_1+\theta_2)
    \end{aligned}
    \label{eq:2dof}
\end{equation}

Esta equivalência permite utilizar posteriormente a resolução geométrica de um manipulador de dois graus de liberdade na cinemática inversa, para determinar os ângulos necessários à posição pretendida de J4.

\subsection{Posição do centro de preensão}

As expressões anteriores permitem calcular a posição das juntas, mas o objetivo da cinemática direta é determinar a posição do centro de preensão da garra.

Para isso, é necessário acrescentar a extensão existente entre J5 e o centro de preensão.

Como os segmentos $L_6$ e $L_7$ estão alinhados com o eixo longitudinal da garra, esta distância é dada por:

\begin{equation}
    D=L_6+L_7
    \label{eq:distancia_garra}
\end{equation}

Considerando inicialmente $q_5=0^\circ$, a extensão da garra permanece no plano lateral, com orientação definida pelo ângulo $\alpha$.

Assim, a posição do centro de preensão é obtida somando as projeções dos elos $L_2$, $L_3$, $L_4$ e da extensão $D$.

As coordenadas resultantes são:

\begin{align}
    d &= L_2\cos\theta_A+
    L_3\cos\theta_B+
    L_4\cos\theta_C+
    D\cos\alpha
    \label{eq:d}\\
    h &= L_1+
    L_2\sin\theta_A+
    L_3\sin\theta_B+
    L_4\sin\theta_C+
    D\sin\alpha
    \label{eq:h}
\end{align}

Na primeira equação, as projeções horizontais dos elos determinam a distância frontal $d$.

Na segunda, somam-se as projeções verticais à altura inicial $L_1$, obtendo-se a altura $h$ do centro de preensão.

A Figura~\ref{fig:soma_projecoes} representa esta decomposição geométrica.

\begin{figure}[htbp]
    \centering
    \resizebox{0.6\textwidth}{!}{\begin{tikzpicture}[scale=0.3, >=stealth,
    elo/.style={line width=2.2pt, line cap=round},
    junta/.style={circle, draw, fill=white, line width=1pt, inner sep=2pt},
    guia/.style={gray!60, thin, densely dotted},
    seg/.style={line width=2.4pt},
    every node/.style={font=\footnotesize}]
  \colorlet{cL1}{gray!70!black}
  \colorlet{cL2}{blue!70!black}
  \colorlet{cL3}{orange!85!black}
  \colorlet{cL4}{green!50!black}
  \colorlet{cD}{violet!80!black}

  % Coordenadas
  \coordinate (B)  at (0,0);
  \coordinate (J2) at (0,13.2);
  \coordinate (J3) at (4.7,23.2);
  \coordinate (J4) at (14.1,21.5);
  \coordinate (J5) at (21.0,24.0);
  \coordinate (G)  at (26.2,9.8);

  % Eixos
  \draw[->, thin] (-2,0) -- (36,0) node[below left] {$d$};
  \draw[->, thin] (B) -- (0,28) node[left] {$h$};
  \fill[gray!25] (-2.5,-0.8) rectangle (2.5,0);

  % Guias
  \foreach \P in {J3,J4,J5,G} \draw[guia] (\P) -- (\P |- 0,-2.2);
  \draw[guia] (J2) -- (29,13.2);
  \draw[guia] (J3) -- (30,23.2);
  \draw[guia] (J4) -- (31,21.5);
  \draw[guia] (J5) -- (32,24.0);
  \draw[guia] (G)  -- (33,9.8);

  % Braços
  \draw[elo, cL1] (B)  -- (J2);
  \draw[elo, cL2] (J2) -- node[left]  {$L_2$} (J3);
  \draw[elo, cL3] (J3) -- node[pos=0.55, above] {$L_3$} (J4);
  \draw[elo, cL4] (J4) -- node[above left] {$L_4$} (J5);
  \draw[elo, cD]  (J5) -- node[right] {$D$} (G);
  \foreach \P in {J2,J3,J4,J5} \node[junta] at (\P) {};
  \fill (G) circle (6pt);
  \node[left=4pt]  at (J2) {J2};
  \node[above=3pt] at (J3) {J3};
  \node[below=4pt] at (J4) {J4};
  \node[above=3pt] at (J5) {J5};
  \node[below left=2pt] at (G) {garra};

  % Ângulos
  \draw[guia] (J2) -- ++(3.5,0);  \draw[->, thin, cL2] ($(J2)+(2.5,0)$) arc[start angle=0, end angle=65, radius=2.5] node[pos=0.4, right] {$\theta_A$};
  \draw[guia] (J3) -- ++(3.5,0);  \draw[->, thin, cL3] ($(J3)+(2.8,0)$) arc[start angle=0, end angle=-10, radius=2.8] node[pos=1, below right=-1pt] {$\theta_B$};
  \draw[guia] (J4) -- ++(3.5,0);  \draw[->, thin, cL4] ($(J4)+(2.8,0)$) arc[start angle=0, end angle=20, radius=2.8] node[pos=0.6, right] {$\theta_C$};
  \draw[guia] (J5) -- ++(3.5,0);  \draw[->, thin, cD]  ($(J5)+(2.5,0)$) arc[start angle=0, end angle=-70, radius=2.5] node[pos=0.6, right] {$\alpha$};

  % Soma de d
  \draw[seg, cL2] (0,-2.2)    -- (4.7,-2.2);
  \draw[seg, cL3] (4.7,-2.2)  -- (14.1,-2.2);
  \draw[seg, cL4] (14.1,-2.2) -- (21.0,-2.2);
  \draw[seg, cD]  (21.0,-2.2) -- (26.2,-2.2);
  \node[cL2, below=3pt] at (2.35,-2.2)  {$L_2\cos\theta_A$};
  \node[cL3, above=3pt] at (9.4,-2.2)   {$L_3\cos\theta_B$};
  \node[cL4, below=3pt] at (17.55,-2.2) {$L_4\cos\theta_C$};
  \node[cD,  above=3pt] at (23.6,-2.2)  {$D\cos\alpha$};
  \draw[|<->|, thin] (0,-5.2) -- node[fill=white, inner sep=1pt] {$d$} (26.2,-5.2);

  % Soma de h
  \draw[seg, cL1, ->] (29,0)    -- (29,13.2);
  \draw[seg, cL2, ->] (30,13.2) -- (30,23.2);
  \draw[seg, cL3, ->] (31,23.2) -- (31,21.5);
  \draw[seg, cL4, ->] (32,21.5) -- (32,24.0);
  \draw[seg, cD,  ->] (33,24.0) -- (33,9.8);
  \draw[guia] (29,13.2) -- (30,13.2) (30,23.2) -- (31,23.2) (31,21.5) -- (32,21.5) (32,24.0) -- (33,24.0);
  \node[cL1, left=1pt]  at (29,6.6)   {$L_1$};
  \node[cL2, right=1pt] at (30,17.5)  {$L_2\sin\theta_A$};
  \node[cL3, right=1pt] at (33.2,22.4) {$L_3\sin\theta_B < 0$};
  \node[cL4, right=1pt] at (32,25.6)  {$L_4\sin\theta_C$};
  \node[cD,  right=1pt] at (33,15.0)  {$D\sin\alpha < 0$};
  \draw[|<->|, thin] (36,0) -- node[fill=white, inner sep=1pt] {$h$} (36,9.8);
  \draw[guia] (33,9.8) -- (36,9.8);
  
\end{tikzpicture}}
    \caption{Decomposição dos elos e soma das projeções no plano lateral.}
    \label{fig:soma_projecoes}
\end{figure}

\subsection{Juntas 5 e 6}

Até este ponto, considerou-se $q_5=0^\circ$, mantendo a extensão da garra no plano lateral. No entanto, quando J5 roda, parte dessa extensão passa a encontrar-se fora do plano.

Nesse caso, a componente de $D$ que permanece no plano lateral é dada por $D\cos q_5$, enquanto a componente perpendicular corresponde a $D\sin q_5$.

Por isso, nas expressões anteriores, a extensão $D$ é substituída pela sua projeção $D\cos q_5$. Obtêm-se, assim:

\begin{align}
    d &= L_2\cos\theta_A+
    L_3\cos\theta_B+
    L_4\cos\theta_C+
    D\cos q_5\cos\alpha
    \label{eq:d_q5}\\
    h &= L_1+
    L_2\sin\theta_A+
    L_3\sin\theta_B+
    L_4\sin\theta_C+
    D\cos q_5\sin\alpha
    \label{eq:h_q5}\\
    l &= L_5-D\sin q_5
    \label{eq:q5}
\end{align}

O parâmetro $l$ representa o desvio lateral total do centro de preensão, considerando o desvio estrutural $L_5$ e a contribuição da rotação de J5.

Na implementação da cinemática inversa desenvolvida neste projeto, considera-se $q_5=0^\circ$. Consequentemente, a garra permanece no plano lateral e o desvio lateral reduz-se a $l=L_5$.

Por sua vez, J6 permite rodar os dedos da garra em torno do eixo longitudinal do \textit{end effector}. Como o centro de preensão se encontra alinhado com este eixo, esta rotação altera apenas a orientação dos dedos, sem modificar a posição do centro de preensão.

Assim, J6 não intervém no cálculo das coordenadas da posição.

